\documentclass[11pt,a4paper]{article}

% --- Encoding and fonts ---
\usepackage[utf8]{inputenc}
\usepackage[T1]{fontenc}

% --- Mathematics ---
\usepackage{amsmath,amssymb,amsthm}
\usepackage{mathtools}

% --- Page layout ---
\usepackage[margin=1in,headheight=14pt]{geometry}

% --- Hyperlinks ---
\usepackage[colorlinks=true,linkcolor=red,citecolor=blue,urlcolor=blue]{hyperref}

% --- Lists ---
\usepackage{enumitem}

% --- Colors and boxes ---
\usepackage{xcolor}
\usepackage[theorems,skins,breakable]{tcolorbox}

% --- Header/footer ---
\usepackage{fancyhdr}
\pagestyle{fancy}
\fancyhf{}
\fancyhead[L]{\textsc{Lesson 9: Convex Optimization and Duality}}
\fancyhead[R]{\thepage}

% --- Tables ---
\usepackage{booktabs}

% ============================================================
% Theorem environments
% ============================================================
\theoremstyle{definition}
\newtheorem{definition}{Definition}[section]
\newtheorem{example}{Example}[section]
\newtheorem{exercise}{Exercise}[section]

\theoremstyle{plain}
\newtheorem{theorem}{Theorem}[section]
\newtheorem{lemma}[theorem]{Lemma}
\newtheorem{proposition}[theorem]{Proposition}
\newtheorem{corollary}[theorem]{Corollary}

\theoremstyle{remark}
\newtheorem{remark}{Remark}[section]

% ============================================================
% Colored boxes
% ============================================================
\newtcolorbox{keyresult}[1][]{
  colback=green!5!white,
  colframe=green!60!black,
  fonttitle=\bfseries,
  title={Key Result},
  breakable,
  #1
}

\newtcolorbox{rlconnection}[1][]{
  colback=blue!5!white,
  colframe=blue!60!black,
  fonttitle=\bfseries,
  title={RL Connection},
  breakable,
  #1
}

\newtcolorbox{warning}[1][]{
  colback=red!5!white,
  colframe=red!60!black,
  fonttitle=\bfseries,
  title={Warning},
  breakable,
  #1
}

% ============================================================
% Custom commands
% ============================================================
\newcommand{\R}{\mathbb{R}}
\newcommand{\N}{\mathbb{N}}
\newcommand{\Q}{\mathbb{Q}}
\newcommand{\Z}{\mathbb{Z}}
\newcommand{\C}{\mathbb{C}}
\newcommand{\Prob}{\mathbb{P}}
\newcommand{\E}{\mathbb{E}}
\newcommand{\indicator}{\mathbf{1}}

\DeclareMathOperator{\dom}{dom}
\DeclareMathOperator{\epi}{epi}
\DeclareMathOperator{\conv}{conv}
\DeclareMathOperator{\prox}{prox}
\DeclareMathOperator*{\argmin}{arg\,min}
\DeclareMathOperator*{\argmax}{arg\,max}
\DeclareMathOperator{\tr}{tr}
\DeclareMathOperator{\diag}{diag}
\DeclareMathOperator{\interior}{int}

% ============================================================
\title{Lesson 30: Interior-Point Methods}
\author{%
  Session 9 of 102 \quad$\mid$\quad Phase I: Mathematical Foundations \quad$\mid$\quad 8h \\[4pt]
  \textit{Primary texts: Boyd \& Vandenberghe Ch.\,4--5, 9--11; Nesterov Ch.\,1--4}
}
\date{}

\begin{document}
\maketitle
\tableofcontents
\newpage

\setcounter{section}{7}

\section{Interior-Point Methods}\label{sec:interior}

\subsection{Logarithmic Barrier}

Interior-point methods solve constrained problems by converting inequality
constraints into a smooth barrier term added to the objective, then following
the ``central path'' as the barrier weight decreases.

\begin{definition}[Logarithmic barrier]\label{def:barrier}
For the convex program~\eqref{eq:std-form} (with equality constraints omitted
for simplicity), the \emph{logarithmic barrier} is
\begin{equation}\label{eq:barrier}
  \phi(x) = -\sum_{i=1}^m \log(-f_i(x)),
\end{equation}
defined for $x$ in the strictly feasible set $\{x : f_i(x) < 0, \, i = 1,\dots,m\}$.
\end{definition}

\begin{proposition}[Barrier is convex]\label{prop:barrier-convex}
If $f_1, \dots, f_m$ are convex, then the logarithmic barrier $\phi$ is convex
on $\{x : f_i(x) < 0\}$.
\end{proposition}

\begin{proof}
The function $-\log(-u)$ is convex and non-decreasing on $u < 0$. Since each
$f_i$ is convex, the composition $-\log(-f_i(x))$ is convex (by the
composition rule for convex non-decreasing outer function and convex inner
function). The sum of convex functions is convex.
\end{proof}

\subsection{Central Path}

\begin{definition}[Barrier problem and central path]\label{def:central-path}
For parameter $t > 0$, the \emph{barrier problem} is
\begin{equation}\label{eq:barrier-prob}
  \min_x \; t\, f_0(x) + \phi(x).
\end{equation}
Assuming a unique minimizer $x^*(t)$ exists for each $t > 0$, the curve
$\{x^*(t) : t > 0\}$ is called the \emph{central path}.
\end{definition}

\begin{theorem}[Central path optimality bound]\label{thm:central-path}
Each point $x^*(t)$ on the central path satisfies
\begin{equation}\label{eq:central-bound}
  f_0(x^*(t)) - p^* \le \frac{m}{t},
\end{equation}
where $p^*$ is the optimal value and $m$ is the number of inequality
constraints.
\end{theorem}

\begin{proof}
At $x^*(t)$, the gradient of the barrier objective vanishes:
\[
  t\, \nabla f_0(x^*(t)) + \nabla \phi(x^*(t)) = 0.
\]
The gradient of the barrier is
$\nabla \phi(x) = \sum_{i=1}^m \frac{-\nabla f_i(x)}{-f_i(x)}
= \sum_{i=1}^m \frac{\nabla f_i(x)}{-f_i(x)}$.

Define the dual variables $\lambda^*_i(t) = \frac{1}{t(-f_i(x^*(t)))} > 0$.
Then the stationarity condition becomes
\[
  \nabla f_0(x^*(t)) + \sum_{i=1}^m \lambda^*_i(t) \nabla f_i(x^*(t)) = 0.
\]
Thus $x^*(t)$ minimizes $L(x, \lambda^*(t), 0) = f_0(x) + \sum_i \lambda^*_i(t) f_i(x)$ over $x$,
so the dual function evaluates to
\begin{align*}
  g(\lambda^*(t), 0)
  &= f_0(x^*(t)) + \sum_i \lambda^*_i(t) f_i(x^*(t)) \\
  &= f_0(x^*(t)) + \sum_i \frac{f_i(x^*(t))}{t(-f_i(x^*(t)))} \\
  &= f_0(x^*(t)) - \frac{m}{t}.
\end{align*}
By weak duality, $g(\lambda^*(t), 0) \le p^*$, so
$f_0(x^*(t)) - m/t \le p^*$, giving $f_0(x^*(t)) - p^* \le m/t$.
\end{proof}

\subsection{Barrier Method}

\begin{definition}[Barrier method]\label{def:barrier-method}
The \emph{barrier method} (or \emph{path-following method}) proceeds as follows:
\begin{enumerate}
  \item Initialize with a strictly feasible $x^{(0)}$, $t^{(0)} > 0$,
    and a growth factor $\sigma > 1$.
  \item At iteration $\ell$: solve (approximately) the barrier problem
    $\min_x \; t^{(\ell)} f_0(x) + \phi(x)$ starting from $x^{(\ell)}$,
    using Newton's method. Set $x^{(\ell+1)}$ to the result.
  \item Update $t^{(\ell+1)} = \sigma \cdot t^{(\ell)}$.
  \item Stop when $m/t^{(\ell)} \le \epsilon$.
\end{enumerate}
\end{definition}

\begin{theorem}[Complexity of the barrier method]\label{thm:barrier-complexity}
To achieve $f_0(x) - p^* \le \epsilon$, the barrier method requires
$O(\sqrt{m}\log(m/\epsilon))$ Newton steps (outer iterations), each involving
the solution of a linear system of size $n \times n$.
\end{theorem}

\begin{proof}[Proof sketch]
By Theorem~\ref{thm:central-path}, we need $t \ge m/\epsilon$, so
$\lceil \log_\sigma(m/(t^{(0)}\epsilon))\rceil$ outer iterations suffice.
The key result from the theory of self-concordant barriers (Nesterov and
Nemirovskii) is that with the optimal choice $\sigma = 1 + 1/\sqrt{m}$,
each centering step (finding $x^*(t)$ to sufficient accuracy) requires
$O(1)$ Newton steps (specifically, a constant number independent of $m$ and
$n$). The total number of outer iterations is then
$\log_\sigma(m/\epsilon) = O(\sqrt{m}\log(m/\epsilon))$,
using $\log(1 + 1/\sqrt{m}) \approx 1/\sqrt{m}$.
\end{proof}

\begin{keyresult}[title={Interior-Point Complexity}]
Interior-point methods solve convex programs with $m$ inequality constraints
to $\epsilon$-accuracy in $O(\sqrt{m}\log(m/\epsilon))$ Newton iterations.
Each Newton iteration costs $O(n^3)$ for a dense $n$-dimensional problem
(or less with structure). This polynomial complexity in $m$, $n$, and
$\log(1/\epsilon)$ makes interior-point methods highly efficient for
moderate-dimensional problems.
\end{keyresult}

\begin{rlconnection}
Interior-point methods are used to solve the linear programs that arise in the
LP formulation of finite MDPs. For large-scale problems, they complement
value iteration and policy iteration. In constrained RL (safe RL), solving
the inner constrained optimization at each policy update step can be handled
by interior-point methods when the constraint set has moderate dimension.
\end{rlconnection}

%% ============================================================
%% SECTION 9: Exercises
%% ============================================================

\section{Exercises}\label{sec:exercises}

\begin{exercise}[Dual of a linear program]\label{ex:lp-dual}
Consider the LP $\min_x \; c^\top x$ subject to $Ax = b$, $x \ge 0$.
\begin{enumerate}[label=(\alph*)]
  \item Write the Lagrangian and derive the dual problem. Show it is
    $\max_\nu \; b^\top \nu$ subject to $A^\top \nu \le c$.
  \item Prove that strong duality always holds for feasible LPs (Hint: LP
    feasible sets are polyhedral, so Slater's condition can be replaced by
    the weaker condition of feasibility).
\end{enumerate}
\end{exercise}

\begin{exercise}[KKT for least-squares with non-negativity]\label{ex:nnls}
Consider the non-negative least-squares problem:
$\min_{x \ge 0} \; \frac{1}{2}\|Ax - b\|^2$.
Write the KKT conditions explicitly and show they reduce to:
$x \ge 0$, $A^\top(Ax - b) \ge 0$, $x_i [A^\top(Ax - b)]_i = 0$ for all $i$.
\end{exercise}

\begin{exercise}[Convergence rate verification]\label{ex:gd-verify}
Let $f(x) = \frac{1}{2}x^\top Q x$ where $Q$ is symmetric positive definite
with eigenvalues $\mu = \lambda_{\min}(Q)$ and $L = \lambda_{\max}(Q)$.
\begin{enumerate}[label=(\alph*)]
  \item Verify that $f$ is $\mu$-strongly convex and $L$-smooth.
  \item Run gradient descent with $\eta = 1/L$ and verify that
    $f(x_k) \le (1 - \mu/L)^k f(x_0)$.
  \item For $Q = \diag(1, 100)$, compute the condition number and determine
    how many iterations are needed for $f(x_k)/f(x_0) \le 10^{-6}$.
\end{enumerate}
\end{exercise}

\begin{exercise}[Soft-thresholding derivation]\label{ex:soft-thresh}
Derive the soft-thresholding operator: show that
$\prox_{\tau|\cdot|}(v) = \text{sign}(v)\max(|v| - \tau, 0)$ for a scalar $v$.
(Hint: consider the three cases $v > \tau$, $|v| \le \tau$, and $v < -\tau$
separately.)
\end{exercise}

\begin{exercise}[Dual of SDP relaxation]\label{ex:sdp-dual}
Consider the boolean optimization problem $\min_{x \in \{-1,+1\}^n} x^\top W x$.
The standard SDP relaxation replaces $x x^\top$ with a matrix variable
$X \succeq 0$, $X_{ii} = 1$.
\begin{enumerate}[label=(\alph*)]
  \item Write this relaxation as a standard SDP.
  \item Derive its Lagrange dual.
  \item Show that weak duality provides a lower bound on the original
    combinatorial problem.
\end{enumerate}
\end{exercise}

\begin{exercise}[Proving the PL inequality]\label{ex:pl}
Let $f$ be $\mu$-strongly convex and differentiable.
\begin{enumerate}[label=(\alph*)]
  \item Prove that $\|\nabla f(x)\|^2 \ge 2\mu(f(x) - f(x^*))$ for all $x$.
  \item Show by example that the PL inequality can hold for non-convex
    functions (consider $f(x) = x^2 + 3\sin^2(x)$ near the origin; verify
    numerically).
  \item Prove that gradient descent converges linearly under the PL inequality
    alone (without assuming convexity).
\end{enumerate}
\end{exercise}

\begin{exercise}[FISTA for LASSO]\label{ex:fista-lasso}
The LASSO problem is $\min_x \frac{1}{2}\|Ax - b\|^2 + \tau\|x\|_1$.
\begin{enumerate}[label=(\alph*)]
  \item Identify $f(x)$ and $h(x)$ in the composite formulation.
  \item Write out the FISTA iteration explicitly (including the
    soft-thresholding step).
  \item What is the Lipschitz constant $L$ of $\nabla f$? Express it in terms
    of $A$.
\end{enumerate}
\end{exercise}

\begin{exercise}[Central path for LP]\label{ex:lp-barrier}
For the LP $\min\; c^\top x$ subject to $x \ge 0$ (i.e., $f_i(x) = -x_i$):
\begin{enumerate}[label=(\alph*)]
  \item Write the barrier function $\phi(x) = -\sum_i \log(x_i)$.
  \item Show that the central path satisfies $x^*_i(t) = 1/(tc_i)$ when $c > 0$.
  \item Verify that $c^\top x^*(t) - p^* = n/t$ where $n$ is the dimension
    and $p^* = 0$.
\end{enumerate}
\end{exercise}

\begin{exercise}[Duality gap in constrained MDP]\label{ex:cmdp}
Consider a constrained MDP with one constraint: $\max_\pi J(\pi)$ subject to
$C(\pi) \le \alpha$, where $J$ and $C$ are the expected return and expected
cost under policy $\pi$ over a finite MDP.
\begin{enumerate}[label=(\alph*)]
  \item Write the Lagrangian $L(\pi, \lambda)$ for this problem.
  \item Show that the dual function $g(\lambda) = \max_\pi L(\pi, \lambda)$
    corresponds to solving an unconstrained MDP with modified reward
    $r(s,a) - \lambda c(s,a)$.
  \item Under what condition does strong duality hold? (Hint: relate to
    Slater's condition in the occupancy-measure LP formulation.)
\end{enumerate}
\end{exercise}

\begin{exercise}[Moreau envelope and gradient]\label{ex:moreau}
The \emph{Moreau envelope} of a convex function $h$ with parameter $\gamma > 0$ is
$M_\gamma h(v) = \min_x \{h(x) + \frac{1}{2\gamma}\|x - v\|^2\}$.
\begin{enumerate}[label=(\alph*)]
  \item Show that $M_\gamma h$ is differentiable even when $h$ is not, with
    gradient $\nabla M_\gamma h(v) = \frac{1}{\gamma}(v - \prox_{\gamma h}(v))$.
  \item Prove that $M_\gamma h$ is $(1/\gamma)$-smooth.
  \item Explain how this connects to the proximal point algorithm:
    $x_{k+1} = \prox_{\gamma h}(x_k)$.
\end{enumerate}
\end{exercise}

%% ============================================================
%% SECTION 10: Summary
%% ============================================================

\section{Summary}\label{sec:summary}

\subsection{Key Theorems Recap}

The central results of this document are:
\begin{enumerate}
  \item \textbf{Weak duality} (Theorem~\ref{thm:weak-duality}): Always
    $d^* \le p^*$.
  \item \textbf{Strong duality} (Theorem~\ref{thm:strong-duality}): Under
    Slater's condition, $d^* = p^*$.
  \item \textbf{KKT conditions} (Theorems~\ref{thm:kkt-necessary}
    and~\ref{thm:kkt-sufficient}): Necessary under strong duality, sufficient
    for convex problems.
  \item \textbf{Gradient descent} (Theorems~\ref{thm:gd-convex}
    and~\ref{thm:gd-strong}): $O(1/k)$ for convex, $O(\rho^k)$ for strongly convex.
  \item \textbf{Nesterov acceleration} (Theorem~\ref{thm:nesterov}): Optimal
    $O(1/k^2)$ for smooth convex.
  \item \textbf{Interior-point methods} (Theorem~\ref{thm:barrier-complexity}):
    $O(\sqrt{m}\log(m/\epsilon))$ Newton steps.
\end{enumerate}

\subsection{Convergence Rate Summary}

\begin{table}[h]
\centering
\caption{Convergence rates for first-order and interior-point methods.}
\label{tab:rates}
\renewcommand{\arraystretch}{1.3}
\begin{tabular}{@{}llll@{}}
\toprule
\textbf{Method} & \textbf{Assumption} & \textbf{Rate} & \textbf{Iterations for $\epsilon$} \\
\midrule
Gradient descent & Convex, $L$-smooth &
  $O(L\|x_0-x^*\|^2/k)$ & $O(1/\epsilon)$ \\
Gradient descent & $\mu$-s.c., $L$-smooth &
  $O((1-\mu/L)^k)$ & $O(\kappa\log(1/\epsilon))$ \\
Nesterov accel. & Convex, $L$-smooth &
  $O(L\|x_0-x^*\|^2/k^2)$ & $O(1/\sqrt{\epsilon})$ \\
Nesterov accel. & $\mu$-s.c., $L$-smooth &
  $O((1-1/\sqrt{\kappa})^k)$ & $O(\sqrt{\kappa}\log(1/\epsilon))$ \\
Prox.\ gradient & Convex, $L$-smooth $+$ $h$ &
  $O(L\|x_0-x^*\|^2/k)$ & $O(1/\epsilon)$ \\
FISTA & Convex, $L$-smooth $+$ $h$ &
  $O(L\|x_0-x^*\|^2/k^2)$ & $O(1/\sqrt{\epsilon})$ \\
Interior point & $m$ ineq.\ constraints &
  $O(\sqrt{m}\log(m/\epsilon))$ & $O(\sqrt{m}\log(1/\epsilon))$ \\
\bottomrule
\end{tabular}
\end{table}

\begin{remark}[Lower bounds]\label{rem:lower}
The $O(1/k^2)$ rate for smooth convex functions and $O((1-1/\sqrt{\kappa})^k)$
for smooth strongly convex functions match the information-theoretic lower bounds
for first-order methods (Nesterov, 1983). These are not improvable by any
method that accesses $f$ only through gradient evaluations.
\end{remark}

\begin{remark}[Stochastic extensions]\label{rem:stochastic}
In ML and RL, we typically have access only to stochastic gradients. SGD with
step size $\eta_k = O(1/k)$ achieves $O(1/\sqrt{k})$ for convex and $O(1/k)$
for strongly convex problems (in expectation). Variance reduction techniques
(SVRG, SARAH) recover the deterministic rates up to logarithmic factors. These
stochastic methods are covered in a subsequent document.
\end{remark}


\end{document}